\chapter{Lernziele}
\begin{todolist}
    \item Ich kann Zahlen von einem gegebenen Zahlensystem in das Dezimalsystem umwandeln. Beispiele:
    \begin{itemize}
        \item $101_2 = 5_{10}$
        \item $204_5 = 54_{10}$
        \item $1A_{16} = 26_{10}$
    \end{itemize}
    \item Ich kann Zahlen vom Dezimalsystem in ein vorgegebenes Zahlensystem umwandeln. Beispiele:
    \begin{itemize}
        \item $5_{10} = 101_2$
        \item $135_{10}=1020_5$
        \item $26_{10}=1A_{16}$
    \end{itemize}
    \item Ich kann Zahlen in einem fremden Zahlensystem (schriftlich) addieren. Zum Beispiel:
    \[123_5 + 431_5 = 1104_5\]
    \item Ich kenne die Einheiten Bits und Bytes und kann berechnen, wie viele unterschiedliche Informationen mit einer gewissen Anzahl Bits codiert werden können.
    \item Ich kann erklären, worum es sich bei \gls{ASCII} und \gls{UTF-8} handelt (Sie müssen die \glsentryshort{ASCII}-Tabelle aber nicht auswendig kennen \emoji{winking-face}).
    \item Ich verstehe Bilder als ein Raster aus Pixeln, die je eine bestimmte Farbe haben, und kenne den Begriff der Farbtiefe.
    \item Ich kann berechnen, wie viele Bits nötig sind, um ein Bild mit einer bestimmten Anzahl möglicher Farben zu speichern.
    \item Ich kann erklären, wie man eine \gls{RGB}-Farbcodierung wie \texttt{\#1243A3} interpretieren kann, und diese einer Farbe zuordnen.
    \item Ich kann die gängigen Grössenordnungen für Datenmengen (KB, MB, GB, TB, \ldots) einordnen.
    \item Ich kann die Wahrheitstabellen der logischen Grundgatter NOT, AND und OR aufstellen.
    \item Ich kann die Symbole der Gatter NOT, AND und OR zeichnen und beschriften.
    \item Ich kann den Halbaddierer aus AND- und XOR-Gattern zusammensetzen und seine Funktion (Summe und Übertrag) erklären.
    \item Ich kann einen Volladdierer aus zwei Halbaddierern und einem OR-Gatter aufbauen und zwei Zahlen mit einem 4-Bit-Addierer Stufe für Stufe addieren (Zwischensumme, Überträge, Ergebnis).
    \item Ich kann die vier Hauptkomponenten der Von-Neumann-Architektur (\gls{CPU}, \gls{RAM}, Massenspeicher, Ein-/Ausgabe) benennen und ihre Aufgaben beschreiben.
    \item Ich kann die Bestandteile der \gls{CPU} (\gls{ALU}, Steuerwerk, Register) erläutern.
    \item Ich kann den Fetch-Decode-Execute-Zyklus mit eigenen Worten beschreiben.
    \item Ich kann den Unterschied zwischen \gls{RAM} und Massenspeicher (\gls{HDD}, \gls{SSD}) bezüglich Geschwindigkeit, Latenz, Flüchtigkeit und Grösse erklären.
    \item Ich kann das Schichtenmodell (Software-Stack) mit den Ebenen Hardware, Treiber, Betriebssystem, \gls{API} und Anwendung skizzieren und erklären.
    \item Ich kann die vier zentralen Aufgaben eines Betriebssystems (Prozessverwaltung, Speicherverwaltung, Dateisystem, Geräteverwaltung) beschreiben.
    \item Ich kann das Prinzip des Multitasking (Zeitscheiben) erläutern.
    \item Ich kann an einem konkreten Beispiel (z.\,B. Tastendruck) das Zusammenspiel von Hardware, Betriebssystem und Anwendung erklären.
\end{todolist}
